\documentclass[11pt,letterpaper]{article}
\newcommand{\labid}{1-H}
\newcommand{\labtitleshort}{FPGA Ingestion \& E2E}
% Shared preamble for Module 1 lab handouts.
% Usage: each handout defines \labid and \labtitleshort, then % Shared preamble for Module 1 lab handouts.
% Usage: each handout defines \labid and \labtitleshort, then % Shared preamble for Module 1 lab handouts.
% Usage: each handout defines \labid and \labtitleshort, then \input{LabHandout_Preamble}
% BEFORE \begin{document}. Compiles with a single pdflatex pass (no glossaries tool).

\usepackage[T1]{fontenc}
\usepackage[margin=1in]{geometry}
\usepackage{titlesec}
\usepackage{enumitem}
\usepackage{booktabs}
\usepackage{tabularx}
\usepackage{graphicx}
\graphicspath{{Images/}{../Images/}}
\usepackage{xcolor}
\usepackage{hyperref}
\usepackage{fancyhdr}
\usepackage{amsmath}
\usepackage{amssymb}
\usepackage{tcolorbox}
\usepackage{caption}
\tcbuselibrary{skins,breakable}

% ── Glossary (per-document, built from shared glossary.tex) ──
% Uses \makenoidxglossaries so no external makeglossaries tool is
% needed: sorting and printing is done by LaTeX itself across two
% pdflatex passes (same cadence as the TOC). Only acronyms that
% are actually \gls{}-used in the handout appear in the printed
% glossary, and each \gls{} reference is a live hyperlink that
% jumps to its entry in the Acronyms section at the end.
\usepackage[acronym,nomain,nopostdot]{glossaries}
\makenoidxglossaries
\input{../../glossary}
% Call \printhandoutglossary just before \end{document} in each handout.
\newcommand{\printhandoutglossary}{%
	\clearpage
	\printnoidxglossary[type=\acronymtype,title={Glossary}]%
}

% ── Colors (match Module 1 lesson plan) ──────────────────────
\definecolor{stevens}{RGB}{163,31,52}
\definecolor{darkgray}{RGB}{60,60,60}
\definecolor{lightgray}{RGB}{240,240,240}
\definecolor{accentblue}{RGB}{30,80,140}

\hypersetup{
	colorlinks=true,
	linkcolor=stevens,
	urlcolor=accentblue,
	citecolor=darkgray,
}

% ── Defaults for header variables (handout overrides these) ──
\providecommand{\labid}{?}
\providecommand{\labtitleshort}{Lab Handout}

% ── Running header / footer ──────────────────────────────────
\pagestyle{fancy}
\fancyhf{}
\fancyhead[L]{\small\textcolor{darkgray}{Module 1 --- Lab Handout \labid}}
\fancyhead[R]{\small\textcolor{darkgray}{\labtitleshort}}
\fancyfoot[C]{\small\thepage}
\renewcommand{\headrulewidth}{0.4pt}

% ── Section formatting ───────────────────────────────────────
\titleformat{\section}{\Large\bfseries\color{stevens}}{\thesection}{0.6em}{}
\titleformat{\subsection}{\large\bfseries\color{darkgray}}{\thesubsection}{0.6em}{}

% ── Box styles ───────────────────────────────────────────────
\tcbset{
	infobox/.style={
		enhanced, colback=blue!3, colframe=accentblue, fonttitle=\bfseries,
		title={#1}, coltitle=white, breakable, left=6pt, right=6pt, top=4pt, bottom=4pt
	},
	warning/.style={
		enhanced, colback=red!3, colframe=stevens, fonttitle=\bfseries,
		title={#1}, coltitle=white, breakable, left=6pt, right=6pt, top=4pt, bottom=4pt
	},
	deliverable/.style={
		enhanced, colback=gray!5, colframe=darkgray, fonttitle=\bfseries,
		title={#1}, coltitle=white, breakable, left=6pt, right=6pt, top=4pt, bottom=4pt
	},
}

% ── Figure helper: always draws a placeholder box (images suppressed) ─
% Usage: \labfigure{filename}{width}{caption}{label}
\newcommand{\labfigure}[4]{%
	\begin{figure}[h!]
		\centering
		\fbox{\begin{minipage}[c][2.2in][c]{#2}
			\centering\color{darkgray}\small
				\textbf{[Figure placeholder]}\\[4pt]
				\texttt{\detokenize{#1}}\\[4pt]
				\textit{#3}
		\end{minipage}}%
		\caption{#3}
		\label{#4}
	\end{figure}%
}

% ── Title banner (call at top of document body) ──────────────
\newcommand{\labheader}[2]{%
	\begin{center}
		{\LARGE\bfseries\color{stevens} Lab Handout #1}\\[3pt]
		{\Large\bfseries #2}\\[6pt]
		{\small\color{darkgray} Capstone --- Module 1: \RF{} Hardware Foundations}
	\end{center}
	\vspace{4pt}
	{\color{darkgray}\rule{\textwidth}{0.4pt}}
	\vspace{10pt}
}

% BEFORE \begin{document}. Compiles with a single pdflatex pass (no glossaries tool).

\usepackage[T1]{fontenc}
\usepackage[margin=1in]{geometry}
\usepackage{titlesec}
\usepackage{enumitem}
\usepackage{booktabs}
\usepackage{tabularx}
\usepackage{graphicx}
\graphicspath{{Images/}{../Images/}}
\usepackage{xcolor}
\usepackage{hyperref}
\usepackage{fancyhdr}
\usepackage{amsmath}
\usepackage{amssymb}
\usepackage{tcolorbox}
\usepackage{caption}
\tcbuselibrary{skins,breakable}

% ── Glossary (per-document, built from shared glossary.tex) ──
% Uses \makenoidxglossaries so no external makeglossaries tool is
% needed: sorting and printing is done by LaTeX itself across two
% pdflatex passes (same cadence as the TOC). Only acronyms that
% are actually \gls{}-used in the handout appear in the printed
% glossary, and each \gls{} reference is a live hyperlink that
% jumps to its entry in the Acronyms section at the end.
\usepackage[acronym,nomain,nopostdot]{glossaries}
\makenoidxglossaries
% Shared acronym macro container for EE800/820 lesson plan modules.
% Usage: type \IDE, \FPGA, \UART, etc. First use is expanded; later uses show acronym only.
% Requires in each document preamble:
%   \usepackage[acronym,toc]{glossaries}
%   \makeglossaries
%   \input{../glossary}
% Print used acronyms near end of document:
%   \printglossary[type=\acronymtype,title={Acronyms Used}]

% Acronym definitions (only used entries appear in the printed glossary).
\newacronym{ask}{ASK}{Amplitude Shift Keying}
\newacronym{bga}{BGA}{Ball Grid Array}
\newacronym{bom}{BOM}{Bill of Materials}
\newacronym{bram}{BRAM}{Block Random Access Memory}
\newacronym{crc}{CRC}{Cyclic Redundancy Check}
\newacronym{css}{CSS}{Chirp Spread Spectrum}
\newacronym{dma}{DMA}{Direct Memory Access}
\newacronym{dsp}{DSP}{Digital Signal Processing}
\newacronym{eda}{EDA}{Exploratory Data Analysis}
\newacronym{fifo}{FIFO}{First In, First Out}
\newacronym{fpga}{FPGA}{Field-Programmable Gate Array}
\newacronym{fsk}{FSK}{Frequency Shift Keying}
\newacronym{fsm}{FSM}{Finite State Machine}
\newacronym{gpio}{GPIO}{General-Purpose Input/Output}
\newacronym{hal}{HAL}{Hardware Abstraction Layer}
\newacronym{i2c}{I\textsuperscript{2}C}{Inter-Integrated Circuit}
\newacronym{ide}{IDE}{Integrated Development Environment}
\newacronym{ila}{ILA}{Integrated Logic Analyzer}
\newacronym{ism}{ISM}{Industrial, Scientific, and Medical}
\newacronym{isr}{ISR}{Interrupt Service Routine}
\newacronym{lut}{LUT}{Look-Up Table}
\newacronym{mcu}{MCU}{Microcontroller Unit}
\newacronym{ml}{ML}{Machine Learning}
\newacronym{ook}{OOK}{On-Off Keying}
\newacronym{per}{PER}{Packet Error Rate}
\newacronym{rf}{RF}{Radio Frequency}
\newacronym{rssi}{RSSI}{Received Signal Strength Indicator}
\newacronym{rtl}{RTL}{Register-Transfer Level}
\newacronym{snr}{SNR}{Signal-to-Noise Ratio}
\newacronym{spi}{SPI}{Serial Peripheral Interface}
\newacronym{uart}{UART}{Universal Asynchronous Receiver/Transmitter}
\newacronym{vhdl}{VHDL}{VHSIC Hardware Description Language}

% Convenience wrappers so you can type \IDE, \RF, etc.
\providecommand{\ASK}{\gls{ask}}
\providecommand{\BGA}{\gls{bga}}
\providecommand{\BOM}{\gls{bom}}
\providecommand{\BRAM}{\gls{bram}}
\providecommand{\CRC}{\gls{crc}}
\providecommand{\CSS}{\gls{css}}
\providecommand{\DMA}{\gls{dma}}
\providecommand{\DSP}{\gls{dsp}}
\providecommand{\EDA}{\gls{eda}}
\providecommand{\FIFO}{\gls{fifo}}
\providecommand{\FPGA}{\gls{fpga}}
\providecommand{\FSK}{\gls{fsk}}
\providecommand{\FSM}{\gls{fsm}}
\providecommand{\GPIO}{\gls{gpio}}
\providecommand{\HAL}{\gls{hal}}
\providecommand{\IIC}{\gls{i2c}}
\providecommand{\IDE}{\gls{ide}}
\providecommand{\ILA}{\gls{ila}}
\providecommand{\ISM}{\gls{ism}}
\providecommand{\ISR}{\gls{isr}}
\providecommand{\LUT}{\gls{lut}}
\providecommand{\MCU}{\gls{mcu}}
\providecommand{\ML}{\gls{ml}}
\providecommand{\OOK}{\gls{ook}}
\providecommand{\PER}{\gls{per}}
\providecommand{\RF}{\gls{rf}}
\providecommand{\RSSI}{\gls{rssi}}
\providecommand{\RTL}{\gls{rtl}}
\providecommand{\SNR}{\gls{snr}}
\providecommand{\SPI}{\gls{spi}}
\providecommand{\UART}{\gls{uart}}
\providecommand{\VHDL}{\gls{vhdl}}


% Call \printhandoutglossary just before \end{document} in each handout.
\newcommand{\printhandoutglossary}{%
	\clearpage
	\printnoidxglossary[type=\acronymtype,title={Glossary}]%
}

% ── Colors (match Module 1 lesson plan) ──────────────────────
\definecolor{stevens}{RGB}{163,31,52}
\definecolor{darkgray}{RGB}{60,60,60}
\definecolor{lightgray}{RGB}{240,240,240}
\definecolor{accentblue}{RGB}{30,80,140}

\hypersetup{
	colorlinks=true,
	linkcolor=stevens,
	urlcolor=accentblue,
	citecolor=darkgray,
}

% ── Defaults for header variables (handout overrides these) ──
\providecommand{\labid}{?}
\providecommand{\labtitleshort}{Lab Handout}

% ── Running header / footer ──────────────────────────────────
\pagestyle{fancy}
\fancyhf{}
\fancyhead[L]{\small\textcolor{darkgray}{Module 1 --- Lab Handout \labid}}
\fancyhead[R]{\small\textcolor{darkgray}{\labtitleshort}}
\fancyfoot[C]{\small\thepage}
\renewcommand{\headrulewidth}{0.4pt}

% ── Section formatting ───────────────────────────────────────
\titleformat{\section}{\Large\bfseries\color{stevens}}{\thesection}{0.6em}{}
\titleformat{\subsection}{\large\bfseries\color{darkgray}}{\thesubsection}{0.6em}{}

% ── Box styles ───────────────────────────────────────────────
\tcbset{
	infobox/.style={
		enhanced, colback=blue!3, colframe=accentblue, fonttitle=\bfseries,
		title={#1}, coltitle=white, breakable, left=6pt, right=6pt, top=4pt, bottom=4pt
	},
	warning/.style={
		enhanced, colback=red!3, colframe=stevens, fonttitle=\bfseries,
		title={#1}, coltitle=white, breakable, left=6pt, right=6pt, top=4pt, bottom=4pt
	},
	deliverable/.style={
		enhanced, colback=gray!5, colframe=darkgray, fonttitle=\bfseries,
		title={#1}, coltitle=white, breakable, left=6pt, right=6pt, top=4pt, bottom=4pt
	},
}

% ── Figure helper: always draws a placeholder box (images suppressed) ─
% Usage: \labfigure{filename}{width}{caption}{label}
\newcommand{\labfigure}[4]{%
	\begin{figure}[h!]
		\centering
		\fbox{\begin{minipage}[c][2.2in][c]{#2}
			\centering\color{darkgray}\small
				\textbf{[Figure placeholder]}\\[4pt]
				\texttt{\detokenize{#1}}\\[4pt]
				\textit{#3}
		\end{minipage}}%
		\caption{#3}
		\label{#4}
	\end{figure}%
}

% ── Title banner (call at top of document body) ──────────────
\newcommand{\labheader}[2]{%
	\begin{center}
		{\LARGE\bfseries\color{stevens} Lab Handout #1}\\[3pt]
		{\Large\bfseries #2}\\[6pt]
		{\small\color{darkgray} Capstone --- Module 1: \RF{} Hardware Foundations}
	\end{center}
	\vspace{4pt}
	{\color{darkgray}\rule{\textwidth}{0.4pt}}
	\vspace{10pt}
}

% BEFORE \begin{document}. Compiles with a single pdflatex pass (no glossaries tool).

\usepackage[T1]{fontenc}
\usepackage[margin=1in]{geometry}
\usepackage{titlesec}
\usepackage{enumitem}
\usepackage{booktabs}
\usepackage{tabularx}
\usepackage{graphicx}
\graphicspath{{Images/}{../Images/}}
\usepackage{xcolor}
\usepackage{hyperref}
\usepackage{fancyhdr}
\usepackage{amsmath}
\usepackage{amssymb}
\usepackage{tcolorbox}
\usepackage{caption}
\tcbuselibrary{skins,breakable}

% ── Glossary (per-document, built from shared glossary.tex) ──
% Uses \makenoidxglossaries so no external makeglossaries tool is
% needed: sorting and printing is done by LaTeX itself across two
% pdflatex passes (same cadence as the TOC). Only acronyms that
% are actually \gls{}-used in the handout appear in the printed
% glossary, and each \gls{} reference is a live hyperlink that
% jumps to its entry in the Acronyms section at the end.
\usepackage[acronym,nomain,nopostdot]{glossaries}
\makenoidxglossaries
% Shared acronym macro container for EE800/820 lesson plan modules.
% Usage: type \IDE, \FPGA, \UART, etc. First use is expanded; later uses show acronym only.
% Requires in each document preamble:
%   \usepackage[acronym,toc]{glossaries}
%   \makeglossaries
%   % Shared acronym macro container for EE800/820 lesson plan modules.
% Usage: type \IDE, \FPGA, \UART, etc. First use is expanded; later uses show acronym only.
% Requires in each document preamble:
%   \usepackage[acronym,toc]{glossaries}
%   \makeglossaries
%   \input{../glossary}
% Print used acronyms near end of document:
%   \printglossary[type=\acronymtype,title={Acronyms Used}]

% Acronym definitions (only used entries appear in the printed glossary).
\newacronym{ask}{ASK}{Amplitude Shift Keying}
\newacronym{bga}{BGA}{Ball Grid Array}
\newacronym{bom}{BOM}{Bill of Materials}
\newacronym{bram}{BRAM}{Block Random Access Memory}
\newacronym{crc}{CRC}{Cyclic Redundancy Check}
\newacronym{css}{CSS}{Chirp Spread Spectrum}
\newacronym{dma}{DMA}{Direct Memory Access}
\newacronym{dsp}{DSP}{Digital Signal Processing}
\newacronym{eda}{EDA}{Exploratory Data Analysis}
\newacronym{fifo}{FIFO}{First In, First Out}
\newacronym{fpga}{FPGA}{Field-Programmable Gate Array}
\newacronym{fsk}{FSK}{Frequency Shift Keying}
\newacronym{fsm}{FSM}{Finite State Machine}
\newacronym{gpio}{GPIO}{General-Purpose Input/Output}
\newacronym{hal}{HAL}{Hardware Abstraction Layer}
\newacronym{i2c}{I\textsuperscript{2}C}{Inter-Integrated Circuit}
\newacronym{ide}{IDE}{Integrated Development Environment}
\newacronym{ila}{ILA}{Integrated Logic Analyzer}
\newacronym{ism}{ISM}{Industrial, Scientific, and Medical}
\newacronym{isr}{ISR}{Interrupt Service Routine}
\newacronym{lut}{LUT}{Look-Up Table}
\newacronym{mcu}{MCU}{Microcontroller Unit}
\newacronym{ml}{ML}{Machine Learning}
\newacronym{ook}{OOK}{On-Off Keying}
\newacronym{per}{PER}{Packet Error Rate}
\newacronym{rf}{RF}{Radio Frequency}
\newacronym{rssi}{RSSI}{Received Signal Strength Indicator}
\newacronym{rtl}{RTL}{Register-Transfer Level}
\newacronym{snr}{SNR}{Signal-to-Noise Ratio}
\newacronym{spi}{SPI}{Serial Peripheral Interface}
\newacronym{uart}{UART}{Universal Asynchronous Receiver/Transmitter}
\newacronym{vhdl}{VHDL}{VHSIC Hardware Description Language}

% Convenience wrappers so you can type \IDE, \RF, etc.
\providecommand{\ASK}{\gls{ask}}
\providecommand{\BGA}{\gls{bga}}
\providecommand{\BOM}{\gls{bom}}
\providecommand{\BRAM}{\gls{bram}}
\providecommand{\CRC}{\gls{crc}}
\providecommand{\CSS}{\gls{css}}
\providecommand{\DMA}{\gls{dma}}
\providecommand{\DSP}{\gls{dsp}}
\providecommand{\EDA}{\gls{eda}}
\providecommand{\FIFO}{\gls{fifo}}
\providecommand{\FPGA}{\gls{fpga}}
\providecommand{\FSK}{\gls{fsk}}
\providecommand{\FSM}{\gls{fsm}}
\providecommand{\GPIO}{\gls{gpio}}
\providecommand{\HAL}{\gls{hal}}
\providecommand{\IIC}{\gls{i2c}}
\providecommand{\IDE}{\gls{ide}}
\providecommand{\ILA}{\gls{ila}}
\providecommand{\ISM}{\gls{ism}}
\providecommand{\ISR}{\gls{isr}}
\providecommand{\LUT}{\gls{lut}}
\providecommand{\MCU}{\gls{mcu}}
\providecommand{\ML}{\gls{ml}}
\providecommand{\OOK}{\gls{ook}}
\providecommand{\PER}{\gls{per}}
\providecommand{\RF}{\gls{rf}}
\providecommand{\RSSI}{\gls{rssi}}
\providecommand{\RTL}{\gls{rtl}}
\providecommand{\SNR}{\gls{snr}}
\providecommand{\SPI}{\gls{spi}}
\providecommand{\UART}{\gls{uart}}
\providecommand{\VHDL}{\gls{vhdl}}


% Print used acronyms near end of document:
%   \printglossary[type=\acronymtype,title={Acronyms Used}]

% Acronym definitions (only used entries appear in the printed glossary).
\newacronym{ask}{ASK}{Amplitude Shift Keying}
\newacronym{bga}{BGA}{Ball Grid Array}
\newacronym{bom}{BOM}{Bill of Materials}
\newacronym{bram}{BRAM}{Block Random Access Memory}
\newacronym{crc}{CRC}{Cyclic Redundancy Check}
\newacronym{css}{CSS}{Chirp Spread Spectrum}
\newacronym{dma}{DMA}{Direct Memory Access}
\newacronym{dsp}{DSP}{Digital Signal Processing}
\newacronym{eda}{EDA}{Exploratory Data Analysis}
\newacronym{fifo}{FIFO}{First In, First Out}
\newacronym{fpga}{FPGA}{Field-Programmable Gate Array}
\newacronym{fsk}{FSK}{Frequency Shift Keying}
\newacronym{fsm}{FSM}{Finite State Machine}
\newacronym{gpio}{GPIO}{General-Purpose Input/Output}
\newacronym{hal}{HAL}{Hardware Abstraction Layer}
\newacronym{i2c}{I\textsuperscript{2}C}{Inter-Integrated Circuit}
\newacronym{ide}{IDE}{Integrated Development Environment}
\newacronym{ila}{ILA}{Integrated Logic Analyzer}
\newacronym{ism}{ISM}{Industrial, Scientific, and Medical}
\newacronym{isr}{ISR}{Interrupt Service Routine}
\newacronym{lut}{LUT}{Look-Up Table}
\newacronym{mcu}{MCU}{Microcontroller Unit}
\newacronym{ml}{ML}{Machine Learning}
\newacronym{ook}{OOK}{On-Off Keying}
\newacronym{per}{PER}{Packet Error Rate}
\newacronym{rf}{RF}{Radio Frequency}
\newacronym{rssi}{RSSI}{Received Signal Strength Indicator}
\newacronym{rtl}{RTL}{Register-Transfer Level}
\newacronym{snr}{SNR}{Signal-to-Noise Ratio}
\newacronym{spi}{SPI}{Serial Peripheral Interface}
\newacronym{uart}{UART}{Universal Asynchronous Receiver/Transmitter}
\newacronym{vhdl}{VHDL}{VHSIC Hardware Description Language}

% Convenience wrappers so you can type \IDE, \RF, etc.
\providecommand{\ASK}{\gls{ask}}
\providecommand{\BGA}{\gls{bga}}
\providecommand{\BOM}{\gls{bom}}
\providecommand{\BRAM}{\gls{bram}}
\providecommand{\CRC}{\gls{crc}}
\providecommand{\CSS}{\gls{css}}
\providecommand{\DMA}{\gls{dma}}
\providecommand{\DSP}{\gls{dsp}}
\providecommand{\EDA}{\gls{eda}}
\providecommand{\FIFO}{\gls{fifo}}
\providecommand{\FPGA}{\gls{fpga}}
\providecommand{\FSK}{\gls{fsk}}
\providecommand{\FSM}{\gls{fsm}}
\providecommand{\GPIO}{\gls{gpio}}
\providecommand{\HAL}{\gls{hal}}
\providecommand{\IIC}{\gls{i2c}}
\providecommand{\IDE}{\gls{ide}}
\providecommand{\ILA}{\gls{ila}}
\providecommand{\ISM}{\gls{ism}}
\providecommand{\ISR}{\gls{isr}}
\providecommand{\LUT}{\gls{lut}}
\providecommand{\MCU}{\gls{mcu}}
\providecommand{\ML}{\gls{ml}}
\providecommand{\OOK}{\gls{ook}}
\providecommand{\PER}{\gls{per}}
\providecommand{\RF}{\gls{rf}}
\providecommand{\RSSI}{\gls{rssi}}
\providecommand{\RTL}{\gls{rtl}}
\providecommand{\SNR}{\gls{snr}}
\providecommand{\SPI}{\gls{spi}}
\providecommand{\UART}{\gls{uart}}
\providecommand{\VHDL}{\gls{vhdl}}


% Call \printhandoutglossary just before \end{document} in each handout.
\newcommand{\printhandoutglossary}{%
	\clearpage
	\printnoidxglossary[type=\acronymtype,title={Glossary}]%
}

% ── Colors (match Module 1 lesson plan) ──────────────────────
\definecolor{stevens}{RGB}{163,31,52}
\definecolor{darkgray}{RGB}{60,60,60}
\definecolor{lightgray}{RGB}{240,240,240}
\definecolor{accentblue}{RGB}{30,80,140}

\hypersetup{
	colorlinks=true,
	linkcolor=stevens,
	urlcolor=accentblue,
	citecolor=darkgray,
}

% ── Defaults for header variables (handout overrides these) ──
\providecommand{\labid}{?}
\providecommand{\labtitleshort}{Lab Handout}

% ── Running header / footer ──────────────────────────────────
\pagestyle{fancy}
\fancyhf{}
\fancyhead[L]{\small\textcolor{darkgray}{Module 1 --- Lab Handout \labid}}
\fancyhead[R]{\small\textcolor{darkgray}{\labtitleshort}}
\fancyfoot[C]{\small\thepage}
\renewcommand{\headrulewidth}{0.4pt}

% ── Section formatting ───────────────────────────────────────
\titleformat{\section}{\Large\bfseries\color{stevens}}{\thesection}{0.6em}{}
\titleformat{\subsection}{\large\bfseries\color{darkgray}}{\thesubsection}{0.6em}{}

% ── Box styles ───────────────────────────────────────────────
\tcbset{
	infobox/.style={
		enhanced, colback=blue!3, colframe=accentblue, fonttitle=\bfseries,
		title={#1}, coltitle=white, breakable, left=6pt, right=6pt, top=4pt, bottom=4pt
	},
	warning/.style={
		enhanced, colback=red!3, colframe=stevens, fonttitle=\bfseries,
		title={#1}, coltitle=white, breakable, left=6pt, right=6pt, top=4pt, bottom=4pt
	},
	deliverable/.style={
		enhanced, colback=gray!5, colframe=darkgray, fonttitle=\bfseries,
		title={#1}, coltitle=white, breakable, left=6pt, right=6pt, top=4pt, bottom=4pt
	},
}

% ── Figure helper: always draws a placeholder box (images suppressed) ─
% Usage: \labfigure{filename}{width}{caption}{label}
\newcommand{\labfigure}[4]{%
	\begin{figure}[h!]
		\centering
		\fbox{\begin{minipage}[c][2.2in][c]{#2}
			\centering\color{darkgray}\small
				\textbf{[Figure placeholder]}\\[4pt]
				\texttt{\detokenize{#1}}\\[4pt]
				\textit{#3}
		\end{minipage}}%
		\caption{#3}
		\label{#4}
	\end{figure}%
}

% ── Title banner (call at top of document body) ──────────────
\newcommand{\labheader}[2]{%
	\begin{center}
		{\LARGE\bfseries\color{stevens} Lab Handout #1}\\[3pt]
		{\Large\bfseries #2}\\[6pt]
		{\small\color{darkgray} Capstone --- Module 1: \RF{} Hardware Foundations}
	\end{center}
	\vspace{4pt}
	{\color{darkgray}\rule{\textwidth}{0.4pt}}
	\vspace{10pt}
}


\begin{document}
\labheader{1-H}{\FPGA{} \UART{} Ingestion and End-to-End Data Flow}

\section{Objective}
Implement a minimal \UART{} ingestion path on the Nexys A7-100T that receives the 43-byte forwarding frames from the STM32 receiver (LH1-G), detects the start delimiter, extracts the beacon \mbox{ID}, and displays it on the seven-segment display. This closes the Module~1 end-to-end data path: \textbf{beacon $\to$ \RF{} link $\to$ receiver $\to$ \FPGA{}}.

\section{Prerequisites}
\begin{itemize}[leftmargin=2em]
	\item Lab Handouts~1-E (Vivado toolchain), 1-F (beacon), and 1-G (receiver) complete.
	\item The receiver node from LH1-G must be producing valid frames on \texttt{USART3 TX} (validated by the Python parser).
\end{itemize}

\section{Architecture}

\begin{table}[h!]
	\centering
	\begin{tabular}{l l}
		\toprule
		\textbf{Block} & \textbf{Function} \\
		\midrule
		\texttt{uart\_rx} & 115200\,8N1 \UART{} receiver (oversampled) \\
		\texttt{byte\_fifo} & Small 16-byte \FIFO{} between \texttt{uart\_rx} and frame detector \\
		\texttt{frame\_detector} & Synchronize on \texttt{0x7E} start, capture 43 bytes, expose beacon \mbox{ID} \\
		\texttt{seg7\_driver} & Drive the four right-most digits with the beacon \mbox{ID} in hex \\
		\bottomrule
	\end{tabular}
\end{table}

\labfigure{LH1H_block_diagram.png}{0.9\textwidth}{\FPGA{} ingestion block diagram: \texttt{uart\_rx} $\to$ \texttt{byte\_fifo} $\to$ \texttt{frame\_detector} $\to$ \texttt{seg7\_driver}. Key signals (\texttt{rx\_strobe}, \texttt{beacon\_id}, \texttt{new\_frame}) are labeled on the data path.}{fig:lh1h-blocks}

\newpage
\section{Wiring}
Connect the receiver's \texttt{USART3 TX} pin to one of the Nexys A7 \mbox{PMOD} input pins. Common ground is essential.
\begin{itemize}[leftmargin=2em]
	\item STM32 \texttt{USART3 TX} (\texttt{PC4}) $\to$ Nexys A7 \texttt{JA[0]} (pin \texttt{C17}, check the \mbox{PMOD} pin mapping in the Nexys A7 reference manual).
	\item STM32 \texttt{GND} $\to$ Nexys A7 \texttt{GND} (\texttt{JA} \mbox{PMOD} pin 5 or 11).
\end{itemize}

\labfigure{LH1H_pmod_wiring.png}{0.8\textwidth}{STM32 \texttt{USART3 TX} (\texttt{PC4}) wired to Nexys A7 \texttt{JA[0]} (pin \texttt{C17}) with common ground. Annotated photo clarifies which \mbox{PMOD} connector is in use.}{fig:lh1h-pmod}

\begin{tcolorbox}[warning={Voltage Levels}]
	The STM32 drives 3.3\,V logic and the Nexys A7 \mbox{PMOD} pins are 3.3\,V-tolerant (in fact, 3.3\,V native). No level shifter is required. Do \textit{not} wire the STM32 to an \mbox{LVCMOS18} bank by mistake --- double-check which \mbox{PMOD} you are using.
\end{tcolorbox}

\section{Procedure}

\subsection{Step 1 --- Project setup}
\begin{enumerate}[leftmargin=2em]
	\item In Vivado, create a new project targeting \texttt{xc7a100tcsg324-1} (or reopen the LH1-E project and add the new sources alongside \texttt{loopback\_top.vhd}).
	\item Add a new top-level entity \texttt{ingest\_top.vhd} with:
	\begin{itemize}[leftmargin=1.2em]
		\item \texttt{clk100} (100\,MHz input from the Nexys A7 oscillator).
		\item \texttt{rst\_n} (active-low reset tied to a pushbutton).
		\item \texttt{uart\_rx} (single input from the \mbox{PMOD} pin).
		\item \texttt{seg} (7 outputs for segments) and \texttt{an} (8 outputs for digit anodes).
	\end{itemize}
	\item Add the pruned \texttt{Nexys-A7-100T-Master.xdc} with the relevant pins uncommented.
\end{enumerate}

\subsection{Step 2 --- \UART{} receiver}
\begin{enumerate}[leftmargin=2em]
	\item Write or reuse a textbook 115200\,8N1 \UART{} receiver. Clock: 100\,MHz, so the baud-tick divisor is $100\,000\,000/115\,200 \approx 868$. Oversample by 16 and sample at bit midpoints.
	\item Expose two outputs: \texttt{rx\_data[7:0]} and \texttt{rx\_strobe} (one-cycle pulse on byte received).
\end{enumerate}

\subsection{Step 3 --- Byte \FIFO{}}
\begin{enumerate}[leftmargin=2em]
	\item Instantiate a small synchronous \FIFO{} (16 entries $\times$ 8 bits is plenty). You can hand-roll it in \VHDL{} or instantiate a Xilinx \mbox{FIFO} primitive --- either is fine for this handout.
	\item Push on \texttt{rx\_strobe}; the frame detector pulls one byte per state tick.
\end{enumerate}

\subsection{Step 4 --- Frame detector \FSM{}}
States: \texttt{IDLE}, \texttt{LEN}, \texttt{PAYLOAD}, \texttt{DONE}.
\begin{enumerate}[leftmargin=2em]
	\item \texttt{IDLE}: pop bytes and discard until a byte equals \texttt{0x7E}. Transition to \texttt{LEN}.
	\item \texttt{LEN}: pop one byte. Expect \texttt{41} (0x29). If it doesn't match, return to \texttt{IDLE}. Otherwise, transition to \texttt{PAYLOAD}.
	\item \texttt{PAYLOAD}: pop 41 more bytes into an internal register bank. The frame layout is: delimiter at 0, length (\texttt{0x29}) at 1, metadata (timestamp, \RSSI{}, \SNR{}, freq.\ error) at 2--9, embedded 32-byte beacon payload at 10--41, frame \CRC{}-8 at 42. The beacon \mbox{ID} is beacon byte 2, i.e.\ frame offset 12. Capture it into a register \texttt{beacon\_id[7:0]}.
	\item \texttt{DONE}: assert a \texttt{new\_frame} pulse for one cycle, then return to \texttt{IDLE}.
\end{enumerate}
For Week~4 scope, ignore the frame \CRC{}-8 and embedded beacon \CRC{}-16 checks inside the \FPGA{}. Module~2 will add both.

\labfigure{LH1H_frame_detector_fsm.png}{0.8\textwidth}{Frame-detector \FSM{}: \texttt{IDLE} $\to$ \texttt{LEN} $\to$ \texttt{PAYLOAD} $\to$ \texttt{DONE}, with transition conditions annotated on each edge.}{fig:lh1h-fsm}

\newpage
\subsection{Step 5 --- Seven-segment driver}
\begin{enumerate}[leftmargin=2em]
	\item Multiplex the four right-most digits at $\sim$1\,kHz per digit (use a clock divider from 100\,MHz).
	\item Display the captured \texttt{beacon\_id} in hex on digits \texttt{[1:0]}. Leave digits \texttt{[3:2]} blank or use them to show a simple frame counter that increments on \texttt{new\_frame}.
\end{enumerate}

\subsection{Step 6 --- Synthesize, implement, program}
\begin{enumerate}[leftmargin=2em]
	\item Follow the same flow as LH1-E: Synthesis, Implementation, Generate Bitstream, Program.
	\item Verify zero timing-constraint violations in the implementation report.
\end{enumerate}

\labfigure{LH1H_timing_report.png}{0.85\textwidth}{Vivado implementation timing summary with positive \mbox{WNS}. Any negative slack must be resolved before programming the board.}{fig:lh1h-timing}

\newpage
\subsection{Step 7 --- End-to-end run}
\begin{enumerate}[leftmargin=2em]
	\item Power on the beacon(s) from LH1-F and the receiver from LH1-G. The STM32 receiver should be \texttt{printf}ing debug output over its virtual COM port.
	\item Power on the Nexys A7 and program the bitstream.
	\item Observe the seven-segment display. It should alternate between \texttt{01} and \texttt{02} as each beacon's frames arrive.
	\item Introduce a deliberate perturbation: turn off Beacon~A. The display should lock to \texttt{02} only. Turn it back on. It should resume alternating.
\end{enumerate}

\labfigure{LH1H_seg7_display.jpg}{0.7\textwidth}{Module~1 capstone shot: Nexys A7 seven-segment display showing a live beacon \mbox{ID} captured from the full \textbf{beacon $\to$ \RF{} link $\to$ receiver $\to$ \FPGA{}} chain.}{fig:lh1h-seg7}

\newpage
\section{Verification Checklist}
\begin{itemize}[leftmargin=2em]
	\item[$\square$] \texttt{uart\_rx} correctly decodes bytes at 115200 baud (verify with ILA or test bench).
	\item[$\square$] Frame detector synchronizes on the \texttt{0x7E} delimiter and captures 43-byte frames without losing alignment.
	\item[$\square$] Seven-segment display shows \texttt{01}/\texttt{02} matching the live beacon \mbox{ID}s.
	\item[$\square$] Removing and restoring a beacon causes the display to respond within 2~seconds.
	\item[$\square$] Vivado implementation report: zero timing violations.
\end{itemize}

\begin{tcolorbox}[warning={Troubleshooting}]
	\begin{itemize}[leftmargin=1.2em]
		\item \textbf{No display change}: verify the \mbox{PMOD} pin is correct in the \texttt{.xdc} and that the STM32 ground is tied to the Nexys A7 ground.
		\item \textbf{Garbage bytes}: baud divisor wrong. Recompute against the actual \texttt{clk100} frequency.
		\item \textbf{Frame detector stuck in \texttt{IDLE}}: length byte mismatch. Confirm the receiver is sending \texttt{0x29} as byte 1.
		\item \textbf{Display flickers unreadably}: multiplex rate too slow (increase per-digit frequency to $\geq$500\,Hz).
	\end{itemize}
\end{tcolorbox}

\begin{tcolorbox}[deliverable={Lab Handout 1-H Deliverable --- and Module 1 Final Deliverable}]
	\begin{enumerate}[leftmargin=1.2em]
		\item All \VHDL{} sources: \texttt{uart\_rx}, \texttt{byte\_fifo}, \texttt{frame\_detector}, \texttt{seg7\_driver}, \texttt{ingest\_top}.
		\item Pruned \texttt{.xdc} constraint file.
		\item Vivado utilization and timing reports.
		\item Short video (30--60\,s) of the complete end-to-end system running: beacons transmitting, receiver forwarding, Nexys A7 seven-segment display updating with beacon \mbox{ID}s.
		\item One-page Module~1 retrospective: what worked, what didn't, what you would change before Module~2.
	\end{enumerate}
\end{tcolorbox}

\printhandoutglossary

\end{document}
